\chapter{Lifting of Quantum Torsion}
\label{chap:torsion}

% ── Jet Prolongation ─────────────────────────────────────────────────────────

\begin{proposition}[Jet Prolongation]
\label{prop:jet-prolongation}
Let $(M, G, E)$ be a $\mathbb{C}$ quantum observable universe
(Definition~\ref{def:c-quantum-observable-universe}) with connection $A$ and
metric $g$ on $E$.  Write $\widetilde{M} = E$ for the total space of $E$,
viewed as a manifold, and define $E'$ as the bundle of orthonormal frames of
the first jet bundle $J^1 E$.  Then $A$ and $g$ uniquely determine a
connection $A'$ on $E'$ such that $(\widetilde{M}, G, E')$ is a $\mathbb{C}$
quantum observable universe.
\end{proposition}

\begin{proof}
% TODO: proof
\end{proof}

% ── Recursive Brane Tower ────────────────────────────────────────────────────

\begin{definition}[Recursive Brane Tower]
\label{def:recursive-brane-tower}
Let $(M, G, E)$ be a $\mathbb{C}$ quantum observable universe
(Definition~\ref{def:c-quantum-observable-universe}).  Set $M_0 = M$ and
$E_1 = E$.  \emph{Inductively}, define $E_{n+1}$ as the jet prolongation
(Proposition~\ref{prop:jet-prolongation}) of $E_n$.

The \emph{brane relation} is the following recursive structure:
\begin{enumerate}
  \item Each total space $E_n$ is treated as the base manifold $M_{n+1}$ for
        the next level; $E_n$ is the \emph{brane} of $E_{n+1}$.
  \item At each level, $A_n \in \mathcal{A}(E_{n-1})$ is the $n$-th level
        instanton connection satisfying the Hermitian Yang--Mills equations on
        $E_n$.
\end{enumerate}
\end{definition}

\begin{example}[First two levels]
\label{ex:brane-tower-levels}
Starting from a $\mathbb{C}$ quantum observable universe $(M, G, E_1)$, the
first jet prolongation yields $E_2$ with base manifold $M_1 = E_1$.  The
connection $A_1$ on $E_1$ satisfies the Hermitian Yang--Mills equations, and
$A_2$ on $E_2$ is uniquely determined by
Proposition~\ref{prop:jet-prolongation}.
\end{example}

% ── Quantum Characteristic Classes ───────────────────────────────────────────

\begin{definition}[Quantum Characteristic Classes]
\label{def:quantum-characteristic-classes}
Let $\{E_n\}_{n \geq 1}$ be the recursive brane tower
(Definition~\ref{def:recursive-brane-tower}).  The \emph{quantum
characteristic classes} are the cohomology classes
\[
  \kappa_n \;\in\; H^*\!\bigl(E_{n+1},\, \mathbb{Z}\bigr)
\]
that label the monodromy fixed points of the holonomy representation as one
transitions from $E_n$ to $E_{n+1}$.
\end{definition}

% ── Betti numbers ────────────────────────────────────────────────────────────

\begin{remark}[Betti numbers of the brane tower]
\label{rem:betti-numbers-brane}
Let $b_k(E_n) = \dim H^k(E_n, \mathbb{R})$ denote the $k$-th Betti number of
$E_n$.  Since $E_{n+1}$ is the bundle of orthonormal frames of $J^1 E_n$, the
Betti numbers of $E_{n+1}$ are determined by those of $E_n$ via the Leray--Serre
spectral sequence applied to the fibration
$E_{n+1} \to E_n$.  In particular, the quantum characteristic classes
(Definition~\ref{def:quantum-characteristic-classes}) obstruct the collapse of
this spectral sequence and control the growth of $b_k(E_{n+1})$ in terms of
$b_k(E_n)$.
\end{remark}

% ── Direct Limit Algebra ─────────────────────────────────────────────────────

\begin{definition}[$\mathbb{C}$ Quantum Observable Universe as Direct Limit]
\label{def:direct-limit-algebra}
Let $\{E_n\}_{n \geq 1}$ be the recursive brane tower
(Definition~\ref{def:recursive-brane-tower}) with substrate parameters $q_n$
(Definition~\ref{def:substrate-parameter}) at each level.  The
\emph{$\mathbb{C}$ quantum observable universe} is the direct limit algebra
\[
  \mathcal{A}_\infty
  \;=\;
  \varinjlim\,
  \Gamma\!\bigl(J^1 E_n,\, \mathrm{SU}_{q_n}(2)\bigr).
\]
\end{definition}

% ── Jet Shell ────────────────────────────────────────────────────────────────

\begin{remark}[Jet Shell of the universe]
\label{rem:jet-shell}
In the limit $n \to \infty$, the infinite jet bundle
$J^\infty E = \varprojlim\, J^k E$ becomes the \emph{jet shell} of the
universe.  All finite-level brane data are encoded in this pro-finite
structure, and the algebra $\mathcal{A}_\infty$
(Definition~\ref{def:direct-limit-algebra}) acts naturally on sections of the
jet shell.
\end{remark}

% ── Vacuum State ─────────────────────────────────────────────────────────────

\begin{remark}[The Vacuum State]
\label{rem:vacuum-state}
The \emph{physical vacuum} is a state
$\omega \colon \mathcal{A}_\infty \to \mathbb{C}$ on the direct limit algebra
(Definition~\ref{def:direct-limit-algebra}) that is invariant under the global
$\mathrm{SU}(2)$ action and the recursive $q$-deformations:
\[
  \omega(g \cdot a) = \omega(a)
  \quad\text{for all } g \in \mathrm{SU}(2),\;
  a \in \mathcal{A}_\infty.
\]
\end{remark}
